\begin{ex}
Tính: $14 \times 4 = \dots$
\shortans{56}
\loigiai{
$14 \times 4 = 56$.
}
\end{ex}

\begin{ex}
Tính: $23 \times 4 = \dots$
\shortans{92}
\loigiai{
$23 \times 4 = 92$.
}
\end{ex}

\begin{ex}
Tính: $68 : 2 = \dots$
\shortans{34}
\loigiai{
$68 : 2 = 34$.
}
\end{ex}

\begin{ex}
Tính: $96 : 3 = \dots$
\shortans{32}
\loigiai{
$96 : 3 = 32$.
}
\end{ex}

\begin{ex}
Tính: $85 : 5 = \dots$
\shortans{17}
\loigiai{
$85 : 5 = 17$.
}
\end{ex}

\begin{ex}
Tính: $72 : 4 = \dots$
\shortans{18}
\loigiai{
$72 : 4 = 18$.
}
\end{ex}

\begin{ex}
Tìm số dư trong phép chia $47 : 6$.
\shortans{5}
\loigiai{
$47 = 6 \times 7 + 5$, dư 5.
}
\end{ex}

\begin{ex}
Tìm số dư trong phép chia $59 : 8$.
\shortans{3}
\loigiai{
$59 = 8 \times 7 + 3$, dư 3.
}
\end{ex}

\begin{ex}
Gấp 15 lên 3 lần được bao nhiêu?
\shortans{45}
\loigiai{
$15 \times 3 = 45$.
}
\end{ex}

\begin{ex}
Giảm 48 đi 4 lần được bao nhiêu?
\shortans{12}
\loigiai{
$48 : 4 = 12$.
}
\end{ex}

\begin{ex}
Tính giá trị biểu thức: $13 \times 5 - 25$.
\shortans{40}
\loigiai{
$65 - 25 = 40$.
}
\end{ex}

\begin{ex}
Tính giá trị biểu thức: $92 : 4 + 37$.
\shortans{60}
\loigiai{
$23 + 37 = 60$.
}
\end{ex}

\begin{ex}
Tính giá trị biểu thức: $(24 + 16) : 5$.
\shortans{8}
\loigiai{
$40 : 5 = 8$.
}
\end{ex}

\begin{ex}
Tính giá trị biểu thức: $60 - 45 : 3$.
\shortans{45}
\loigiai{
$60 - 15 = 45$.
}
\end{ex}

\begin{ex}
Một cuộn dây dài 64 m cắt thành 4 đoạn bằng nhau. Hỏi mỗi đoạn dài bao nhiêu mét?
\shortans{16}
\loigiai{
$64 : 4 = 16$ m.
}
\end{ex}

\begin{ex}
Mẹ mua 14 quả cam, mẹ mua số quả quýt gấp 3 lần số quả cam. Hỏi mẹ mua bao nhiêu quả quýt?
\shortans{42}
\loigiai{
$14 \times 3 = 42$ quả quýt.
}
\end{ex}

\begin{ex}
Một cửa hàng có 90 kg gạo, đã bán đi một phần ba số gạo đó. Hỏi cửa hàng đã bán bao nhiêu ki-lô-gam gạo?
\shortans{30}
\loigiai{
$90 : 3 = 30$ kg.
}
\end{ex}

\begin{ex}
Lớp 3C có 35 học sinh chia đều thành 5 nhóm. Hỏi mỗi nhóm có bao nhiêu học sinh?
\shortans{7}
\loigiai{
$35 : 5 = 7$ học sinh.
}
\end{ex}

\begin{ex}
Một thùng đựng 54 lít mật ong, người ta rót đều vào 6 can. Hỏi mỗi can có bao nhiêu lít mật ong?
\shortans{9}
\loigiai{
$54 : 6 = 9$ lít.
}
\end{ex}

\begin{ex}
Trong phép chia có số chia là 7, thương là 9 và số dư là 4. Số bị chia là bao nhiêu?
\shortans{67}
\loigiai{
$9 \times 7 + 4 = 67$.
}
\end{ex}

\begin{ex}
Có 28 học sinh nam và 20 học sinh nữ cùng xếp thành 6 hàng đều nhau. Mỗi hàng có bao nhiêu học sinh?
\shortans{8}
\loigiai{
$(28 + 20) : 6 = 48 : 6 = 8$ học sinh.
}
\end{ex}

\begin{ex}
Một hình chữ nhật có chiều rộng 7 cm, chiều dài gấp 2 lần chiều rộng. Chu vi hình chữ nhật đó là bao nhiêu xăng-ti-mét?
\shortans{42}
\loigiai{
Chiều dài = 14 cm. Chu vi = $(14 + 7) \times 2 = 42$ cm.
}
\end{ex}

\begin{ex}
Tìm $x$: $x \times 4 = 96$.
\shortans{24}
\loigiai{
$x = 96 : 4 = 24$.
}
\end{ex}

\begin{ex}
Tìm $x$: $x : 3 = 28$.
\shortans{84}
\loigiai{
$x = 28 \times 3 = 84$.
}
\end{ex}

\begin{ex}
Tính: $15 \times 4 + 15 = \dots$
\shortans{75}
\loigiai{
$60 + 15 = 75$.
}
\end{ex}
